\clearpage\fancyhead[L]{\small KOMET.}\SourceLeaf{350}
\begin{center}\textbf{KOMET.}\par\vspace{4pt}\hrule\vspace{4pt}\fontsize{10}{12}\selectfont Monatliche Correspondenz. Band XXV, S. 94..97, Januar 1812.\par\vspace{4pt}\hrule\vspace{4pt}\end{center}

\par\indent [Der Inhalt dieser Briefe ist bis auf den letzten Absatz auch im Artikel vom 25. Januar 1812 der Gött. gel. Anz. wiedergegeben. Hier steht jedoch als Zeitpunkt des Periheldurchgangs .\;.\;.\;.\;. 1811 Nov. 11,392108, Meridian von Göttingen.]

\par\indent Herr \textsc{Wachter} hat vor Kurzem ebenfalls die Bedeckung von $\lambda$ Aquarii berechnet und die folgenden verbesserten Konjunktionszeiten gefunden.

\begin{center}
\fontsize{10}{12}\selectfont\setlength{\tabcolsep}{6pt}\renewcommand{\arraystretch}{1.15}\begin{tabular}{l|ccc}
Göttingen & $10^{\text{u}}$ & $22^{\text{m}}$ & $29^{\text{s}}69$ \\
Mannheim & 10 & 16 & 38.01 \\
Cronach & 10 & 28 & 3.50 \\
\end{tabular}
\end{center}

\noindent Die ersten beiden stimmen gut mit Ihren Ergebnissen überein, bei der letzten gibt es aber einen Unterschied von 10 Sekunden. Da er die Rechnung zweimal mit zwei verschiedenen Methoden durchgeführt hat, finden Sie es vielleicht der Mühe wert, Ihre letzte Rechnung noch einmal zu prüfen.
